\section{Part V: Intellectual Property, Technology Transfer, and Translation}

\subsection{36. An invention is not a product, a patent is not a moat}

A scientific result may be protectable but never marketable; a software tool may sell without patent protection; a patent may be narrow, expensive to maintain, or vulnerable to substitutes. Technology transfer means moving usable knowledge into practice under appropriate rights and incentives. Its success can be licensing, open-source adoption, published standards, industry deployment, a spinout, or improved public capability-not only royalties.

Ask what economic or social problem the invention solves, what complementary assets deployment requires, who has the incentive to invest, and what IP approach best accelerates use.

\subsection{37. Background and foreground IP}

\textbf{Background IP} is knowledge, software, datasets, tools, and rights each party brings to a project. \textbf{Foreground IP} arises from the contracted work, subject to legal and contract rules. Collaboration agreements must specify both, along with improvements, derivative works, moral rights where applicable, open-source obligations, patent prosecution, publication, confidentiality, and licensing rights.

If a sponsor buys exclusive ownership of everything developed, the fee should reflect the value and opportunity cost of that transfer. If the institution promises open research while granting secret perpetual exclusivity, it has created incompatible duties.

\subsection{38. Student and staff inventors}

Research institutions should not reflexively claim ownership of every student project or of ideas created by visitors. Rights can depend on employment, facilities, agreements, funding, invention circumstances, and applicable law. Students must have clear prior disclosure, access to independent advice when significant equity or IP is involved, and freedom from grading coercion. Staff inventors need fair attribution, defined rewards, and transparent disclosure obligations.

WIPO's IP-policy materials specifically address ownership, 3rd-party engagement, commercialization, incentives, accounting, and conflicts [R4]. Treat them as design guidance, not a universal legal template.

\subsection{39. Publication timing versus patent timing}

Disclosing an invention publicly before evaluating protection can destroy patent possibilities in some jurisdictions. Delaying every result for speculative IP can damage scholarship. Implement a modest invention-disclosure process: researcher reports potentially protectable subject matter before conference or code release; a qualified reviewer screens patentability and commercial relevance rapidly; the team decides to file, keep trade secret where ethical and lawful, license, or publish openly.

Define default short review periods and escalation paths. Scientists must know what they may publish and when. Inventorship and authorship are not interchangeable legal or scholarly concepts.

\subsection{40. Technology-readiness levels need economic complements}

Engineering readiness asks whether a technology works in increasingly realistic environments. Commercial readiness additionally requires a buyer, integration, reliability, service, regulation, manufacturing or deployment capability, and positive unit economics. A validated laboratory prototype often sits far from a repeatable enterprise sale. A more sophisticated demonstration may increase costs enough to reduce commercial appeal.

Build a \textbf{translation ledger} with 4 columns: scientific evidence, engineering robustness, adoption and willingness to pay, and finance/operations. Do not mark a project ``commercialization-ready'' because the paper was accepted.

\subsection{41. Licensing architecture}

Licenses can be nonexclusive, exclusive by field, exclusive by territory, or time-limited with performance obligations. Terms may include upfront fees, milestones, running royalties, minimum royalties, sublicensing rights, audit rights, development milestones, improvement rights, reserved academic-use rights, and reversion on non-use. An exclusive license that a partner shelves can be worse than a lower-priced nonexclusive license that diffuses the technology.

A technology-transfer decision should model expected social and financial impact, enforcement costs, licensee execution capacity, and bargaining power. The optimal arrangement differs across pharmaceutical inventions, open technical standards, educational assessments, and applied software.

\subsection{42. The TTO as a function, not a premature department}

At the 1st handful of inventions, a named responsible person with outside legal counsel and a clear disclosure workflow may suffice. A large Technology Transfer Office with lawyers, patent managers, licensing officers, and spinout support becomes rational only when the volume and complexity of work justify it. Overbuilding the office can create pressure to patent trivial ideas to justify staffing.

The function needs an invention ledger, contract repository, reporting cadence, conflict register, patent-cost budget, and documented criteria for abandon, publish, protect, or license decisions.

\subsection{43. Spinout versus license versus internal product}

\textbf{Spinout} when a distinct high-growth product requires a dedicated entrepreneurial team, external capital, and governance separate from academic work. \textbf{License} when another company possesses manufacturing, distribution, regulatory, or customer relationships that the institution lacks. \textbf{Internal product} when it strategically strengthens existing operations and can be managed without undermining research. \textbf{Open release} when diffusion, standards, public mission, or community contribution dominates appropriation.

No single route is inherently prestigious. A spinout that raises funding but never solves a problem is not a better result than a widely adopted open method.

\subsection{44. Researcher incentives and conflicts}

The research lead may hold equity in a licensee, consult for a sponsor, recruit students into a venture, and serve on a selection committee. These roles can conflict. Require written disclosure, independent review of related transactions, recusal where appropriate, student protections, clear time commitments, and publication and resource-use rules. A founder should not treat a research trainee as free labor for a private company.

Academic credit and compensation mechanisms should reward reproducibility, mentoring, credible contribution, and appropriate translation-not just patents filed or money raised.
